\begin{abstract}
The software around a language model in an agent is called a scaffold, a framework, a runtime or a harness, and none of these names comes with a definition of what the thing contains. We take the position that an LLM agent is a model wrapped in a harness, $\textit{Agent} = \textit{Harness}(\textit{Model})$, and that every harness is built from five mechanistic primitives: \prim{Input}, which captures what comes in; \prim{Context}, which assembles what the model is shown; \prim{Model interface}, which sends the request and brings back the response; \prim{Output}, which carries the response out; and \prim{Control flow}, which decides whether to go again, hand back or stop. Each primitive is defined by what it does and by what it never does, so that the five are mutually exclusive. We then separate three levels at which harnesses can be described: the \emph{primitives}, which every harness has; the \emph{design decisions} made within each primitive, which is where harnesses differ; and the \emph{production concerns} (persistence, safety and security, observability, resilience, performance, evaluation), which are questions asked of every primitive rather than parts of the harness. As evidence, we deconstruct six harnesses from their primary sources at pinned versions: Claude Code, Codex CLI, opencode, pi, Hermes Agent, and a LangGraph workflow. Every component we found has a home in one primitive or one concern, and the differences between the six are design decisions within the primitives. The vocabulary gives a procedure for locating an agent's failure in one of five places before attributing it to the model, and a schema for reporting the harness alongside the model. The minimal harnesses from which the definitions were drawn are reproduced in an appendix, and the evidence for every assignment is published with the paper.
\end{abstract}
